\section{Testing and Evaluation}

This chapter presents the testing methodology adopted to validate the system’s functionality and robustness. 
It analyzes the observed performance results and discusses the inherent limitations of the current implementation, 
along with potential directions for future improvement.

\subsection{Test Scenarios}

To ensure correctness and reliability, a series of manual, end-to-end test scenarios were executed. 
These tests, derived from procedures documented in \texttt{README\_P2P.md}, cover all major workflows—from user authentication and peer discovery to failure recovery and connection timeouts. 
The primary test cases are summarized in Table~\ref{tab:test-scenarios}.

\begin{longtable}{|p{0.8cm}|p{2.4cm}|p{4cm}|p{6.4cm}|}

\caption{Primary system test scenarios.}
\label{tab:test-scenarios} \\

\hline
\textbf{ID} & \textbf{Scenario} & \textbf{Description} & \textbf{Expected Result} \\ 
\hline
\endfirsthead

\caption[]{-- continued from previous page} \\
\hline
\textbf{ID} & \textbf{Scenario} & \textbf{Description} & \textbf{Expected Result} \\ 
\hline
\endhead

\hline 
\multicolumn{4}{r}{{Continued on next page}} \\ 
\endfoot

\hline
\endlastfoot

% The table body starts here
TS1 & Login and Peer Discovery & A user logs in with valid credentials while another user is already online. & The new user is authenticated successfully, and their presence is broadcast to all active peers, appearing in the peer list in real time. \\ 
\hline
TS2 & P2P Connection Handshake & User~A requests a connection to User~B, and User~B accepts. & User~A’s request appears in User~B’s interface; upon acceptance, a direct TCP connection is established and both UIs transition to the active chat state. \\ 
\hline
TS3 & P2P Connection Rejection & User~A requests a connection to User~B, and User~B rejects. & User~A receives an immediate rejection notification, and the pending chat window closes automatically. \\ 
\hline
TS4 & Real-Time Messaging & During an active P2P session, User~A sends a message to User~B. & The message appears almost instantly in User~B’s chat window, confirming correct operation of the TCP channel. \\ 
\hline
TS5 & Graceful Session Termination & User~A clicks the "End Session" button during an active chat. & The connection closes gracefully; User~B receives a "Session ended by peer" message and the input field is disabled. \\ 
\hline
TS6 & Unexpected Client Disconnect (Heartbeat Timeout) & User~A closes their browser tab without logging out. & After 45~seconds of inactivity, the backend Tracker marks User~A as offline and broadcasts a \texttt{peer-left} event to all other peers. \\ 
\hline
TS7 & P2P Network Interruption (Keepalive Timeout) & The TCP connection between two daemons is forcibly terminated (simulated by killing one daemon process). & After 30~seconds without \texttt{PONG} responses, the remaining daemon detects the idle connection, closes its socket, and updates the UI accordingly. \\ 

\end{longtable}

\subsection{Result Analysis}

The executed tests confirmed that the system fulfills all functional requirements and demonstrates robust behavior under both normal and failure conditions.

\paragraph{Functional Compliance}
All mandatory features were successfully verified, including cookie-based authentication, hybrid Client–Server and P2P communication, 
and real-time messaging without the use of prohibited frameworks or protocols such as WebSocket. 
Each workflow behaved as expected across multiple user sessions.

\paragraph{Latency and Response Time}
The system’s perceived responsiveness is primarily determined by the frontend’s long-polling intervals. 
As defined in \texttt{chat.js}, peer discovery events have a maximum latency of approximately 3~seconds, while chat message retrieval incurs a maximum delay of about 2~seconds. 
Although not instantaneous, this performance is acceptable for the project’s scope and represents a significant improvement over conventional short-polling mechanisms.

\paragraph{System Stability}
The dual-layer health check strategy proved highly effective in maintaining operational stability. 
The 45-second HTTP heartbeat reliably detected and cleaned up inactive (“zombie”) sessions, 
while the 30-second TCP keepalive prevented indefinite hanging of broken peer connections. 
The system remained stable for extended periods with multiple concurrent users and exhibited no runtime crashes or deadlocks.

\paragraph{Scalability and Resource Consumption}
While functionally complete, the current architecture has inherent scalability limitations. 
All application state—including peer registries and message queues—is stored in the backend’s memory, 
creating a centralized bottleneck that constrains concurrent user capacity. 
Moreover, each authenticated user spawns a dedicated \texttt{P2PDaemon} process with multiple threads, 
causing CPU and memory usage to increase linearly with the number of active sessions. 
This design remains suitable for small-scale experiments but would require redesign for deployment at larger scale.

\subsection{Limitations and Future Improvements}

Although the system meets all defined objectives, several limitations were identified during evaluation, 
offering clear directions for enhancement in future iterations.

\paragraph{Current Limitations}
\begin{itemize}
    \item \textbf{No NAT Traversal:} Direct TCP communication functions correctly on a local network but fails across public networks due to Network Address Translation (NAT) restrictions.
    \item \textbf{No Data Persistence:} All user data, session state, and message history are stored in memory and lost when the server restarts.
    \item \textbf{Polling-Induced Latency:} The long-polling approach inherently introduces 2–3~second delays, limiting responsiveness compared to real-time protocols.
\end{itemize}

\paragraph{Future Improvements}
\begin{itemize}
    \item \textbf{Implement NAT Traversal:} Incorporate UDP hole punching or integrate STUN/TURN servers to support peer connectivity across the public internet.
    \item \textbf{Introduce Persistence:} Replace in-memory state with a lightweight database (e.g., SQLite) or an in-memory data store (e.g., Redis) to preserve chat history and session data across restarts.
    \item \textbf{Adopt Real-Time Protocols:} If permitted, migrate from long-polling to WebSocket or HTTP/2 push to eliminate latency and reduce overhead.
    \item \textbf{Enhance Scalability:} Re-architect the monolithic backend into a distributed system. 
    Use shared caching for global state management and consider asynchronous I/O models (e.g., \texttt{asyncio}) for handling higher concurrency with lower resource consumption.
\end{itemize}

In summary, the testing results demonstrate that the system performs reliably and satisfies all project requirements. 
Although constrained by its design scope, the implementation provides a solid foundation for a more scalable and production-ready architecture in future development.
